\chapter{知识的获取概述}
\label{chap:knowledge-acquisition-overview}

假设我们要训练一个识别当地花卉物种的模型。使用时，把一张花卉照片交给模型，希望它返回“鸢尾”或“玫瑰”等物种名称。照片就是这项任务的预测输入，物种名称就是希望得到的预测输出。现在我们已经拍下了一批照片，但这些文件只有像素，没有附上物种答案。如果模型把一张鸢尾照片判成玫瑰，学习算法还缺少一项依据，无法按物种是否正确来检查这次预测。

一个直接的办法是请植物学家看照片，把物种名称写在每张照片旁边。这样，照片甲就由单独的图像变成了“照片甲及其物种答案”。训练时，学习算法把模型的预测与这个答案比较，再据此调整模型内部的参数。参数是决定模型怎样从输入计算输出的数值；训练就是利用样本调整这些数值的过程。这里已经有两个不同的问题：怎样根据答案训练模型，以及怎样取得值得使用的答案。本章讨论后一个问题。

任务要求改变时，缺少的依据也会改变。整张照片的物种答案不能告诉模型花朵位于哪个区域；若希望模型从不同角度认出同一株花，还要说明哪些照片拍的是同一对象。有时没有人能逐张给出物种名称，但我们可以遮住照片的一部分，让模型根据其余部分恢复原内容。类别答案、对象关系、区域位置和恢复要求都可以影响学习，不过它们提供的信息不同，必须按各自用途取得和检查。

本章把为学习补入这些依据的过程称为知识获取。这里的“知识”指与任务相联系、可以用于训练的答案、关系、约束或目标，并不意味着凡是写入训练文件的内容都已经正确。我们需要说明补入了什么、依据来自哪里、在哪些对象上适用，以及经过了什么检查。沿着这批花卉照片，先弄清需要保存和区分的内容，再讨论如何选择来源、组织获取和验收；最后用代表性研究说明这些办法怎样扩展、为什么仍有使用条件。阅读本章只需理解函数输入输出、向量、概率和按条件选择方案的基本思想，各类方法的完整算法由后续章节展开。

\Needspace{6\baselineskip}
\section{知识获取的基本概念}
\label{sec:knowledge-acquisition-problem}

拍摄花卉得到的颜色、纹理和形状是我们实际取得的观测。花卉的物种、所在地区和光照条件描述被研究的环境；镜头取景、曝光和传感器成像决定我们能观察到什么；选择哪些花、在哪些时刻拍照，则决定这批样本怎样取得。例如，只在晴天拍路边常见花，即使照片很多，也没有观察到雨天或稀有物种。把环境、观测方式和选样方式分开，才能判断缺少的是照片，还是照片所对应的任务依据。相关分工可参考\BookChapterRef{chap:mathematics-and-future-research}{卷2《数学准备：概率、因果与决策》的“数学基础与后续研究”章}。

\subsection{从观测到训练依据}
\label{sec:knowledge-acquisition-observation-supervision}

先看最直接的类别答案。给照片甲写入“鸢尾”，得到的是\term{标签}（label），即这张照片在指定分类任务中的目标类别。为了利用它，学习算法需要给预测的偏差计分，这个计分函数称为损失函数（loss function）。一个用于说明思想的计分办法是：预测类别与标签相同计0，不同计1；多张照片上的平均分就表示错分比例。实际训练可以采用其他便于调整参数的损失函数，但都要先规定预测与目标怎样比较。标签提供比较依据，损失函数规定比较方法，两者承担不同职责。

也可能有两张照片乙、丙，我们不知道花的物种，却知道它们拍的是同一株花。希望模型认出同一对象时，可以要求它对两张照片产生相近的描述。模型用一组数值描述图像，这组数值称为表示（representation），可写成特征向量；\term{表示学习}（representation learning）就是通过训练取得这种描述方式，而非完全由人预先规定每个特征。设两张照片的表示分别为$\bm{h}_{\text{乙}},\bm{h}_{\text{丙}}\in\mathbb{R}^{d}$，其中$d$是特征数。为了惩罚同一对象的表示差异，可以在训练中加入
\[
  \ell_{\mathrm{pair}}
  =\lVert\bm{h}_{\text{乙}}-\bm{h}_{\text{丙}}\rVert_2^2.
\]
这里$\lVert\cdot\rVert_2$计算向量的欧氏长度，作用于两组表示之差，就得到它们的欧氏距离；$\ell_{\mathrm{pair}}$是这对照片的距离平方。例如，两组二维表示为$(1,2)$与$(2,2)$时，距离平方为1；调整为相同向量时，这项损失为0。配对关系告诉算法哪些表示应当接近，并没有提供物种名称。还须确认配对确实指向同一对象，而且不能只最小化这一项：把所有照片都映射到同一个向量也能使它为0，却失去了区分不同花卉的能力。完整训练还需保留区分其他对象的信息。

另一类依据规定哪些预测允许同时成立。如果任务已经规定一株花只能属于一个物种，“玫瑰”和“鸢尾”就不能同时成为该对象的最终类别。这种限制是预测约束（prediction constraint），它排除不合法的答案组合，却没有直接告诉我们该选哪一个类别。约束针对的是一株花；一张同时拍到玫瑰和鸢尾的照片可以包含两种花，不能把对象层面的互斥条件误用到整张照片上。采用不同任务或不同对象单位时，约束也要重新说明。

没有物种标签时，照片原有内容还可以提供另一种比较依据。我们先保留完整照片，再遮住其中一块花瓣，把剩余图像交给模型，要求它预测被遮住的内容。原来保留的图像块就是恢复目标（reconstruction target），预测图像块与它的差异可以用于计算损失。只有遮挡后的图像进入预测输入，原图像块另作目标，才能形成需要根据上下文完成的任务。这属于\term{自监督学习}（self-supervised learning）：从观测自身构造学习目标，减少对逐例人工答案的依赖。恢复颜色和纹理可能帮助学习图像表示，但恢复得好还不能证明模型已经会辨认物种。

这些内容进入训练、规定模型应当回答什么或满足什么要求时，都承担\term{监督信号}（supervision signal）的作用。这里的“监督”不限于人工提供物种答案；对象关系、预测约束和自动构造的恢复目标也能指导参数调整。尚未确认适用性与正确性的监督信号称为\term{候选监督信号}（candidate supervision signal）。能够据某个标签计算损失，只说明训练可以使用这个答案，并不说明答案可靠；错误标签也可以让算法执行一次参数更新。标签怎样进入学习的通用分工可参考\BookChapterRef{chap:machine-learning-paradigms}{卷3《基础、理论与可信性》的“范式简介”章}。

\subsection{训练依据的保存与状态}
\label{sec:knowledge-acquisition-record-status}

取得答案以后，我们还要让另一位研究者能够正确读取和复核它。只把“鸢尾”写进文件，会丢掉几个必要信息：它属于照片甲还是照片乙，依据哪份物种名录，由谁在什么材料下作答，适用于哪个地区。如果名录版本改变，同一个名称也可能对应不同的分类范围。因此，训练依据需要连同这些说明一起保存；本部分把这样的记录称为\term{样本标注}（sample annotation）。这个用法比单个类别标签宽：一条记录也可以保存两张照片之间的关系或一个恢复目标。

以人工给照片甲作答为例，一条可追溯的记录可以写成：对象为照片甲，内容为“鸢尾”，来源为标注者甲及其所见照片和所用名录，适用范围为本地区、名录版本$v$下的物种分类。作答者已提交答案，说明获取操作完成、内容已经给出；在独立复核之前，这个答案仍是候选。对象、内容和范围供训练解释答案，来源供复核追查依据，状态则说明当前取得和检查到了哪一步。下面的记号只是把这几项合写成一条记录。

\begin{definition}[样本标注]
  \label{def:knowledge-acquisition-record}
  本部分的样本标注保存一项训练依据及其来源、范围和状态，记为
  \[
    r=(o,s,c,\Omega,z_{\mathrm{sem}},z_{\mathrm{ver}}).
  \]
  其中$o$是所针对的对象，$c$是答案、关系、约束或目标等内容，$\Omega$规定内容适用的样本范围与任务范围。

  来源字段$s=(s_{\mathrm{src}},s_{\mathrm{cond}},z_{\mathrm{act}},z_{\mathrm{exec}})$依次保存来源标识、取得条件、来源动作和执行状态。

  语义状态$z_{\mathrm{sem}}$说明任务内容是否确定、对象是否适用，取值为已给出、未知或不适用；验证状态$z_{\mathrm{ver}}$说明指定检查的进展，取值为候选、已核验或已拒绝。
\end{definition}

在刚才的例子中，$o$对应照片甲，$c$对应“鸢尾”，$\Omega$对应地区、名录版本和任务范围。$s_{\mathrm{src}}$保存标注者身份，$s_{\mathrm{cond}}$保存他看到的材料与作答条件；$z_{\mathrm{act}}$为“作答”，$z_{\mathrm{exec}}$为“成功”。最后两个字段分别为$z_{\mathrm{sem}}=\text{已给出}$与$z_{\mathrm{ver}}=\text{候选}$。这里的“语义”指内容在任务中的含义，“验证”指按指定依据进行的检查。它们分开以后，系统就能保存一个已经提交、但尚未核验的答案。

为了看清为什么需要多个状态，假设照片甲的关键花瓣被遮住。植物学家认真检查了照片，仍不能辨认物种，于是主动不提交类别判断。他完成了请求，因此执行状态是“成功”；他没有给出类别，因此内容是\term{未知}（unknown），即任务答案尚未确定；主动不作判断的动作则称为\term{弃权}（abstention）。执行成功与物种未知可以同时成立。如果程序根本没有打开照片就失败，虽然也没有答案，却没有发生检查后的主动弃权，后续应先修程序或重试。

再假设误把一张石头照片送进只标注花卉的任务。石头不在任务对象范围内，应记为\term{不适用}（not applicable）；继续寻找它的花卉物种没有意义。未知对象可能需要补拍局部或转交专家，不适用对象需要排除或转交另一任务，程序失败需要恢复执行。把它们都存成一个空值，会让后续处理无法选择行动。表~\ref{tab:knowledge-acquisition-record-states}先列各层状态，再用照片甲的几种情形说明它们怎样组合，后续章节沿用这套约定。

\begin{table*}[htbp]
  \centering
  \small
  \setlength{\tabcolsep}{4pt}
  \begin{tabularx}{\linewidth}{@{}>{\RaggedRight\setlength{\parindent}{0pt}\arraybackslash}p{25mm}>{\RaggedRight\setlength{\parindent}{0pt}\arraybackslash}p{55mm}>{\RaggedRight\setlength{\parindent}{0pt}\arraybackslash}X@{}}
    \toprule
    \textbf{样本标注层} & \textbf{允许取值} & \textbf{回答的问题} \\
    \midrule
    执行状态 & 未请求、未处理、处理中、成功、失败 & 请求是否发起；发起后是否已开始、正在处理、完成或失败 \\
    来源动作 & 作答、弃权；无动作时为空 & 来源提交了判断，还是检查后主动不提交判断 \\
    语义状态 & 已给出、未知、不适用 & 内容是否已确定，或者对象本就不属于任务范围 \\
    验证状态 & 候选、已核验、已拒绝 & 针对指定属性的检查是否完成，结果是否允许使用 \\
    规则触发结果 & 触发、未触发、无法确定 & 规则条件是否成立；以具名结果保存在$s_{\mathrm{cond}}$中 \\
    \bottomrule
  \end{tabularx}

  \medskip

  \begin{tabularx}{\linewidth}{@{}>{\RaggedRight\setlength{\parindent}{0pt}\arraybackslash}p{30mm}>{\RaggedRight\setlength{\parindent}{0pt}\arraybackslash}p{19mm}>{\RaggedRight\setlength{\parindent}{0pt}\arraybackslash}p{17mm}>{\RaggedRight\setlength{\parindent}{0pt}\arraybackslash}p{30mm}>{\RaggedRight\setlength{\parindent}{0pt}\arraybackslash}p{18mm}>{\RaggedRight\setlength{\parindent}{0pt}\arraybackslash}X@{}}
    \toprule
    \textbf{样本标注情形} & \textbf{执行} & \textbf{动作} & \textbf{内容／语义} & \textbf{验证} & \textbf{条件或核验依据} \\
    \midrule
    尚未向人工发起请求 & 未请求 & — & 空／未知 & 候选 & 尚未发起任务 \\
    人工提交物种判断 & 成功 & 作答 & 鸢尾／已给出 & 候选 & 依据照片与名录作答 \\
    人工因遮挡而弃权 & 成功 & 弃权 & 空／未知 & 候选 & 辨认物种所需的花瓣局部不可见 \\
    规则运行但未触发 & 成功 & — & 空／未知 & 候选 & 条件检查结果：未触发 \\
    复核确认第二条正确 & 成功 & 作答 & 鸢尾／已给出 & 已核验 & 独立证据支持第二行的物种判断 \\
    \bottomrule
  \end{tabularx}
  \caption{分层保存获取与核验状态。下表都以照片甲为对象，采用同一名录和地区范围；前四行是独立情形，末行表示第二行经过核验后的状态。执行成功可以与内容未知同时成立。破折号表示动作字段为空，末行保留原作答动作，复核依据另行保存。}
  \label{tab:knowledge-acquisition-record-states}
\end{table*}

表中的第四行说明，完成一次条件检查与提交一个物种判断不是同一动作。规则条件不成立时记“未触发”，无需把空答案伪装成弃权；对象若在适用范围外，则根据任务前提记为不适用。核验同样要注明所检查的属性。复核者确认植物学家的弃权符合规范，只能确认动作合规，不能把未知的物种改成已知，也不能把这项检查当成物种正确性的核验。

多份来源作答合在一起后，如果还不能形成可用结论，我们称结果为“未决”，即综合判断尚未确定；这不是单个来源的一次弃权。审计是否开始、是否完成，与“正确、错误、无法裁决”等质量结论也要另行保存，具体处理见第~\ref{chap:knowledge-quality-assessment-correction}章。类似地，“观测不足”“规范歧义”说明没有答案的原因，“不是玫瑰”却已经给出了一项否定判断。如果把这些内容与未知都写成类别0，学习算法就无法分辨它读到的是负标签，还是根本没有任务答案。

例如，把“是否玫瑰”作为一个二分类问题时，标签0可以明确表示“不是玫瑰”，仍然是一项可用于训练的答案。对物种无法判断的照片，应保留未知并说明缺少什么依据，不能仅为方便存储也编码成0。否则，学习算法会把证据不足的照片当成已经确认的负例，原本没有答案的情形就被变成了错误监督。

有了样本标注与状态的区分，再看图~\ref{fig:knowledge-acquisition-observation-to-knowledge}中的三层对象：上层像素是观测，中层答案、关系、约束与目标是候选样本标注，下层才是训练对样本标注的使用。沿四条分支向下看，同一批照片可以支持不同训练要求，样本标注形式不能由文件位置来判断。中层虚线框表示仍待核验：能够自动构造恢复目标，只说明有了一项可训练的要求，不能据此认定它正确或有助于物种识别。

\begin{figure*}[htbp]
  \centering
  \begin{tikzpicture}[
  x=1mm,y=1mm,
  every node/.style={font=\figurefont\footnotesize,text=BookInk},
  layer label/.style={anchor=west,font=\figurefont\bfseries\footnotesize,text=BookTeal},
  observation/.style={rectangle,draw=FigureInput,line width=0.9pt,fill=FigureInput!9,minimum width=52mm,minimum height=12mm,align=center},
  candidate/.style={rectangle,rounded corners=1.2mm,draw=BookAmber,dashed,line width=0.8pt,fill=BookAmber!7,text width=30mm,minimum height=20mm,align=center,inner sep=2mm},
  use/.style={rectangle,draw=FigureOutput,line width=0.8pt,fill=FigureOutput!6,text width=30mm,minimum height=16mm,align=center,inner sep=2mm},
  flow/.style={knowledge arrow,draw=BookMuted,line width=0.8pt,rounded corners=1.2mm}
]
  \node[layer label] at (0,7) {原始观测};
  \node[observation] (photos) at (88,3) {一组花卉照片\\像素与拍摄条件};

  \draw[BookRule,line width=0.6pt] (0,-9) -- (164,-9);
  \node[layer label] at (0,-18) {候选监督信号};
  \node[candidate] (answer) at (30,-27) {类别答案\\“照片甲是鸢尾”\\{\scriptsize 语义：已给出；验证：候选}};
  \node[candidate] (relation) at (70,-27) {对象关系\\“乙、丙来自同一株”\\{\scriptsize 语义：已给出；验证：候选}};
  \node[candidate] (constraint) at (110,-27) {类别约束\\“玫瑰与鸢尾互斥”\\{\scriptsize 语义：已给出；验证：候选}};
  \node[candidate] (restore) at (150,-27) {恢复目标\\“预测被遮蔽花瓣”\\{\scriptsize 语义：已给出；验证：候选}};

  \draw[flow] (photos.south) -- ++(0,-5) -| (answer.north);
  \draw[flow] (photos.south) -- ++(0,-5) -| (relation.north);
  \draw[flow] (photos.south) -- ++(0,-5) -| (constraint.north);
  \draw[flow] (photos.south) -- ++(0,-5) -| (restore.north);

  \draw[BookRule,line width=0.6pt] (0,-40) -- (164,-40);
  \node[layer label] at (0,-49) {训练使用};
  \node[use] (class loss) at (30,-57) {监督分类损失\\学习物种答案};
  \node[use] (relation loss) at (70,-57) {关系一致性目标\\拉近同一对象表示};
  \node[use] (constrained prediction) at (110,-57) {受约束预测\\排除不合法组合};
  \node[use] (reconstruction) at (150,-57) {内容恢复目标\\学习图像表示};

  \draw[flow,draw=BookTeal] (answer.south) -- (class loss.north);
  \draw[flow,draw=BookTeal] (relation.south) -- (relation loss.north);
  \draw[flow,draw=BookTeal] (constraint.south) -- (constrained prediction.north);
  \draw[flow,draw=BookTeal] (restore.south) -- (reconstruction.north);
\end{tikzpicture}

  \caption{观测、候选样本标注与训练使用的三个层次。由同一批花卉照片得到的四类样本标注，分别连接到分类、关系一致性、受约束预测和内容恢复；中层分开标明语义状态与验证状态，虚线框表示样本标注仍需核验。向下的连接说明使用方式，不保证识别收益。内容均为教学构造。}
  \label{fig:knowledge-acquisition-observation-to-knowledge}
\end{figure*}

\subsection{从任务要求确定获取缺口}
\label{sec:knowledge-acquisition-gap-diagnosis}

现在可以把“还不够用”落实到具体缺口。如果目标地区连照片都没有，缺少的是原始观测，需要先采集花卉照片；本章重点讨论已有观测缺少什么解释依据、关系、约束或目标，以及现有标注样本的覆盖和质量哪里不足。一个稀有物种没有任何已确认样本标注，是内容缺失；有了样本标注却只有一次拍摄、没有覆盖特殊光照，是覆盖不足；有整图类别却没有花朵位置，是粒度不足。这三种情况要求补充的内容不同，不能靠重复拍摄同一盆常见花一并解决。

这里的覆盖（coverage）考察目标范围中的哪些情形已经有依据，例如哪些物种、光照和拍摄角度已有标注。粒度（granularity）考察一条依据细到什么对象层次，例如整张照片、一朵花的区域，还是某片花瓣的位置。即使100张照片已经覆盖所有目标物种，若每张只有整图类别，仍不能作为花朵位置的直接答案。反过来，给同一盆玫瑰精确标出100个区域，也没有补入稀有鸢尾的依据。覆盖与粒度是两个不同的检查方向。

有些缺口发生在已经取得的答案中：相似物种被混淆，是内容错误；分类名录调整以后仍沿用旧名称，则是时效失配。拍摄设备或目标地区改变，也可能使旧样本标注的适用条件需要重查，不能只复制旧标签便认为更新完成。表~\ref{tab:knowledge-acquisition-gap-types}把五类缺口接到相应获取需求；验收一列说明每种修补需要观察什么，而不是预先指定某类来源最优。这些缺口可以同时出现，例如稀有物种既缺少多种光照下的样本标注，现有的唯一答案又可能错误。

\begin{table*}[htbp]
  \centering
  \small
  \setlength{\tabcolsep}{3pt}
  \begin{tabularx}{\linewidth}{@{}>{\RaggedRight\setlength{\parindent}{0pt}\arraybackslash}p{25mm}>{\RaggedRight\setlength{\parindent}{0pt}\arraybackslash}p{43mm}>{\RaggedRight\setlength{\parindent}{0pt}\arraybackslash}p{45mm}>{\RaggedRight\setlength{\parindent}{0pt}\arraybackslash}X@{}}
    \toprule
    \textbf{不足类型} & \textbf{花卉识别中的表现} & \textbf{需要补充或改变的内容} & \textbf{可观察的验收依据} \\
    \midrule
    内容缺失 & 目标地区的某个稀有物种没有任何已确认样本标注 & 取得该物种的类别依据、适用范围和来源身份与作答条件 & 目标类别出现已核验样本标注，且对象与范围字段完整 \\
    覆盖不足 & 常见物种有大量近重复照片，稀有物种或特殊光照几乎未覆盖 & 补入困难物种、环境和外观变化的样本标注，而非重复已有高频情形 & 预先划定的物种与拍摄条件各组达到约定的已标注样本数量要求 \\
    粒度不足 & 只有整图类别，无法回答花朵位置或同一对象关系 & 补入与目标问题匹配的区域、实例或成对关系 & 留出的区域或关系样本能按预先规定的对象与语义复核 \\
    内容错误 & 相似物种被混淆，或自动答案把背景叶片当成花朵 & 复核并修正内容，保留原来源、错误类型和修订轨迹 & 在相同独立参考样本上核对修正前后答案，确认已知错误得到处理 \\
    时效失配 & 分类体系已调整，旧样本标注仍用废弃名称；部署地区的物种分布也已变化 & 更新类别映射、适用范围，或重新获取受影响样本标注 & 旧类别均映射到新体系或明确废弃；目标时间窗口的抽查满足预设标准 \\
    \bottomrule
  \end{tabularx}
  \caption{从监督信号缺口确定获取需求。每行由花卉任务中的不足，连接到应补入或改变的内容及验收依据；缺少稀有物种样本标注与缺少位置不能用同一种补充代替。已标注样本数、复核语义和抽查标准须随目标预先确定，各类缺口并非互斥。}
  \label{tab:knowledge-acquisition-gap-types}
\end{table*}

\subsection{从可读取到可靠、有用}
\label{sec:knowledge-acquisition-acceptance}

补入样本标注后，可以按三个不同的问题验收。先检查记录能否读取：照片标识是否存在，标签是否属于约定的类别，名录版本和来源是否齐全，状态组合是否合法。这称为\term{形式完整性}（formal completeness），即记录满足预先规定的字段与格式要求。一条各字段都填齐的“照片甲是玫瑰”记录仍可能把鸢尾认错，因此格式检查之后还需要内容检查。

检查类别内容时，可以请未参与原标注的专家，依据另行提供的清晰局部照片和同一名录重新判断，再比较两份答案。这称为\term{参考一致性}（reference consistency），即候选内容与指定参考证据是否相符。可靠性（reliability）则要求说明这种检查支持了哪些质量判断：例如独立抽查中有多少类别答案满足约定的正确性标准。参考证据的选择也有限制；若两位专家都只看到同一张模糊照片，或者都沿用了同一个错误名录，答案一致仍可能是共享错误。因而“已核验”总要注明检查的属性、依据和标准，不能解释为所有方面永远正确。

即使一批标签按参考证据正确，也还要检查它们是否帮助当前学习任务。假设已有许多清晰玫瑰照片，新标签全部来自相似的玫瑰照片，可能只增加了重复材料，对稀有鸢尾识别的帮助很小。把交付标注放入既定学习算法，再检查训练后模型在目标任务中的结果，评价的是\term{下游效用}（downstream utility）。这里“下游”指取得训练依据之后的模型学习和实际预测；效用衡量这些依据在该任务中带来的收益。

评价这种收益时，要区分训练样本、验证样本和测试样本。训练样本直接参与参数调整；验证样本用于选择获取来源、筛选方案或训练设置；测试样本在方案确定以后，用于检查面对新对象的结果。三者应按用途隔离，并说明如何抽取。随机拆开照片文件仍可能把同一株花的近重复照片分到两侧，使评价高估对新对象的识别能力；因此需要检查对象是否重叠、拍摄条件是否符合未来任务。这些名字本身不保证样本彼此独立或具有相同分布。

形式完整性、参考一致性和下游效用因而逐步回答“能否读取”“指定证据是否支持”“是否帮助任务”。它们检查的对象不同：第一项检查记录，第二项检查内容与依据的关系，第三项检查使用这些记录训练后的模型。训练损失下降也不能代替后两项，因为模型可能只是更贴近已经写入的错误答案。

这几类检查也适用于类别以外的监督。对象关系要核对两张照片是否确实属于同一株花；恢复目标要核对遮挡位置与原内容是否正确对应，再检查这种训练是否服务于识别。模型给出的\term{置信度}（confidence）是它对某个候选答案的确信程度或所赋的分数，具体含义取决于模型及评分方式。高分可以用于决定优先检查哪些对象，却不能自行证明答案正确；模型也可能对训练中反复出现的错误给出高分。

\subsection{在资源限制下选择获取方案}
\label{sec:knowledge-acquisition-plan-selection}

现在假设人工和已有模型都能为稀有物种补入候选标签，但经费、工时和计算资源有限。选择方案不能只比较生成速度，还要比较最后留下的标注是否可靠、是否覆盖缺失情形，以及用于训练有什么收益。一个获取方案应当写明询问哪些对象、调用哪些来源、怎样取得答案、怎样复核，以及什么内容允许交付。仅写“采用模型标注”还不足以确定完整方案。

对每种方案，先区分产生的全部候选与验收后保留的标注。例如，模型提出1000条标签，其中600条因不适用或证据不足而没有通过交付要求，交付集合只有剩下的400条。获取1000条候选和检查被排除的600条都耗费了资源。若只把调用费用除以1000，就无法说明取得一条可交付标注的完整成本。

为了得到可比较的效用，需要固定学习算法、评价样本和比较基线。这里的基线是约定的参照方案，例如使用已有标签训练的模型。对花卉分类，可用留出准确率相对基线增加了多少来计效用；从80\%提高到85\%，增加的是5个百分点。也可以使用平均损失的相反数，让损失降低对应效用升高。可靠性要明确检查类别正确性、关系正确性还是目标构造是否合规，覆盖则要预先划定物种与拍摄条件组，以及每组需要多少标注。不同任务可以采用不同度量，但比较同一任务的方案时应保持口径一致。

一个教学比较可以说明为什么要把可靠性和覆盖单列。假设两个方案均满足预算，交付标准要求独立抽查的标签正确率至少为95\%，四个预设物种组每组至少有20条合格标注。方案甲的抽查正确率为98\%，训练后的留出准确率为88\%，却遗漏了一个稀有物种组；方案乙的抽查正确率为96\%，留出准确率为85\%，四组均达到数量要求。即使甲的总体准确率更高，它仍未完成指定覆盖要求，应先排除，再比较达标方案。这些数值只演示选择规则；有限抽查存在统计不确定性，实际方案还需规定抽查方式及接受标准。

这些量确定后，我们选择预算以内、可靠性和覆盖均达标的方案，再按任务效用从高到低比较，因此目标使用最大化。可靠性与覆盖单独设为约束，是为了避免总体效用较高的方案仍遗漏稀有物种，或交付不可接受的错误；它们不能仅靠效用平均值来代替。下面用统一符号表示这一选择问题。可行方案（feasible plan）在这里指人员、模型、规则和材料允许执行的完整安排，预算和交付门槛再进一步限制其中可以选用哪些。

\begin{definition}[监督信号获取目标]
  \label{def:knowledge-acquisition-objective}
  给定目标任务$T$、可行获取方案集合$\mathcal{A}$和资源上限$\bm{B}$。每个$a\in\mathcal{A}$指定来源、对象、获取动作与复核方式，产生候选样本标注集合$\mathcal{K}_{\mathrm{cand}}(a)$及按交付要求选出的子集$\mathcal{K}(a)\subseteq\mathcal{K}_{\mathrm{cand}}(a)$。

  设$U_T(\mathcal{K}(a))$为交付样本标注的任务效用，$R(\mathcal{K}(a))$为其可靠性度量，$G(\mathcal{K}(a))$为其目标范围覆盖度量，三者均按数值越大越好的方向定义；$r_0$与$g_0$分别为最低可靠性和覆盖要求。获取目标写为
  \begin{equation}
    \begin{aligned}
      \max_{a\in\mathcal{A}}\quad
      & U_T\bigl(\mathcal{K}(a)\bigr) \\
      \text{s.t.}\quad
      & \bm{C}(a)\preceq\bm{B}, \\
      & R\bigl(\mathcal{K}(a)\bigr)\geq r_0,\qquad
        G\bigl(\mathcal{K}(a)\bigr)\geq g_0,
    \end{aligned}
    \label{eq:knowledge-acquisition-objective}
  \end{equation}
  其中$\bm{C}(a)=(C_{\mathrm{dev}}(a),C_{\mathrm{acq}}(a),C_{\mathrm{review}}(a),C_{\mathrm{maint}}(a))$依次计量准备开发、逐项获取、复核修正和持续维护的成本，$\preceq$表示各项都不超过$\bm{B}$中的对应上限。成本按完整方案核算，包括后来被拒绝的候选所耗费的获取与复核投入。人工工时、调用费用与计算资源应分别报告；同一成本项涉及不同单位时也须列明，除非已有可审计的换算规则，否则不能直接相加。
\end{definition}

式~\eqref{eq:knowledge-acquisition-objective}因而把“获得多少条”改为“交付样本标注在任务中有什么收益、是否达标、耗费多少”。这个目标说明如何比较方案，不保证存在达标方案或最大值能够取到，也不提供寻找最优方案的算法；实际效用和可靠性还需要通过试点、抽查与留出评价估计。固定预算时，可以比较达标方案取得的效用，同时报告覆盖和可靠性；固定目标效果时，则反过来比较达到相同要求需要多少分项投入。

读这个式子时，可以先检查三条限制，再看第一行的最大化。$a$代表完整方案，$\mathcal{A}$是当前人员、模型、规则和材料条件下允许考虑的方案集合；$\mathcal{K}_{\mathrm{cand}}(a)$保存全部候选，$\mathcal{K}(a)$保存按要求保留的部分。$R$与$G$分别检查可靠性和覆盖是否达到$r_0$与$g_0$；只有通过这些检查并满足预算的方案，才进入效用$U_T$的比较。式中的$\text{s.t.}$表示“满足以下约束”。这个选择顺序也说明，效用再高都不能补偿一个必须满足却未满足的条件。

成本向量用来保留不同投入，而不是把不同单位强行合成一个数。例如，开发规则需要20工时、调用模型需要100元，20与100没有可以直接相加的共同单位；复核还可能同时耗费工时和调用费用。$\bm{C}(a)\preceq\bm{B}$表示每一项都分别不超过相应上限，工时有剩余不能自动抵消调用费用超限。准备开发、逐项获取、复核修正和持续维护的分工，则帮助我们看出成本发生在一次性投入还是每轮重复操作中。

批量获取之前的小规模试点，用来确认上述比较能否实际执行：目标样本范围是否清楚，哪些未知和错误可以接受，困难物种或条件是否优先覆盖，以及达到什么条件应当停止。当前批次验收后可以交付，但这只完成了当前范围内的任务。分类体系和样本分布继续变化时，长期方案还要规定何时重新抽查、更新或废弃旧样本标注，并把这部分投入计入持续维护成本。

\section{知识获取的研究内容}
\label{sec:knowledge-acquisition-research-scope}

假设试点已经确认稀有物种缺少类别依据，也缺少不同角度下的对象关系，下一步就要决定从哪里补入这些内容。一条样本标注的判断依据来自哪里，称为它的\term{监督来源}（supervision source）。我们可以询问人工、执行规则、让已有模型预测，或利用照片本身构造目标；取得标注以后，再判断哪些能共同使用、哪些需要修正。来源负责提出内容，融合负责组织多份内容或已有模型，质量检查负责判断是否满足要求。把这几种职责分开，就能看清一个完整方案究竟改变了哪一步。

\subsection{按判断依据选择监督来源}
\label{sec:knowledge-acquisition-source-selection}

人工来源让人根据材料和任务规范作出判断。例如，植物学家观察照片中的花瓣、叶片和生长位置，查阅名录后提交“鸢尾”及依据；必要时也可以要求补拍局部，或者因证据不足而弃权。取得答案依赖人的观察与专业判断，因而需要安排可回答的问题、所见材料和人员条件。人工有机会处理规则尚未写明的边界情形，但不会因为作答者是人就自动正确，逐例判断与复核也持续耗费时间。

规则来源把已经规定的条件反复应用到对象。假设照片附有采集编号，另有一份经确认的“采集编号到物种名称”名录，一条规则可以按编号查表，匹配成功时给出候选类别，找不到时不给类别判断。它重复执行的是已写明的关联条件，而非每次重新作植物学判断。是否误配编号、名录是否过期、条件是否适用于当前照片，都影响结果。规则还可以给出前节的类别互斥约束；它提供哪种内容，取决于所编码的条件和结论。

模型来源把已经学得的预测能力用于新对象。例如，已有视觉模型读取照片，计算各物种的分数，再提出分数最高的类别作为候选标签；它也可能给出候选花朵区域。这里的依据是已有模型从过去训练中学得的规律，不是这张照片上已经核验的物种答案。把这类预测转成新训练样本的目标标签，称为\term{伪标签}（pseudo-label），即由模型提出、用作标签的预测结果。它可以减少逐例人工作答，但若过去的训练只覆盖清晰正面照片，新照片中的遮挡仍可能导致高置信错误。

样本本身也可以提供依据。前节把原图像块保留为恢复目标，不需要先取得物种标签；若拍摄日志明确记录照片乙、丙来自同一株花，这个关系也能用于要求两组表示接近。使用的是原内容或已有关系，而非另一个模型对物种的猜测。不过，构造所依赖的关系仍要成立。随意把外观相似的两朵花配成同一对象，可能把不同物种拉近；为构造不同视图而裁掉全部花朵，也可能只留下背景。目标容易构造与目标适合识别，是两项不同判断。

四类来源按内容主要通过什么机制取得来区分。同一个“鸢尾”标签可以由人工、查表规则或视觉模型提供；一段程序既可能执行规则，也可能遮挡图片来构造自监督目标，不能只凭“用了程序”判定来源。构造过程调用模型时，还应说明模型是提出任务答案，还是仅帮助准备材料。表~\ref{tab:knowledge-acquisition-chapter-map}把这些判断接到后续章节，表~\ref{tab:knowledge-acquisition-source-conditions}再把准备投入、易错条件与检查办法放到同一任务下比较。

\begin{table*}[htbp]
  \centering
  \small
  \setlength{\tabcolsep}{4pt}
  \begin{tabularx}{\linewidth}{@{}>{\RaggedRight\setlength{\parindent}{0pt}\arraybackslash}p{35mm}>{\RaggedRight\setlength{\parindent}{0pt}\arraybackslash}X>{\RaggedRight\setlength{\parindent}{0pt}\arraybackslash}p{22mm}@{}}
    \toprule
    \textbf{研究内容} & \textbf{核心问题与主要内容} & \textbf{后续章节} \\
    \midrule
    从人工获取监督信号 & 怎样把需求变成可回答的任务，并选择交互方式、样本、标注者与投入 & 第~\ref{chap:knowledge-acquisition-human}章 \\
    从规则获取监督信号 & 怎样把模式、名录关联和约束写成可重复执行的条件，并检查候选判断 & 第~\ref{chap:rule-based-knowledge-acquisition}章 \\
    从模型获取监督信号 & 怎样准备已有模型、组织输入和请求、取得候选答案，再根据检查结果调整获取过程 & 第~\ref{chap:knowledge-acquisition-model-based}章 \\
    从样本本身获取监督信号 & 怎样用原内容、已有关系或分组依据构造训练目标，并检查构造条件 & 第~\ref{chap:data-derived-knowledge-acquisition}章 \\
    融合多来源标注与模型 & 怎样使多份标注指向相同对象和含义，再综合使用；已有模型又怎样共同预测、组合计算或提供新模型的训练依据 & 第~\ref{chap:knowledge-fusion}章 \\
    评估与修正监督信号 & 怎样建立判断依据、估计质量、定位问题，并验证修正是否有效 & 第~\ref{chap:knowledge-quality-assessment-correction}章 \\
    \bottomrule
  \end{tabularx}
  \caption{由获取需求查找后续章节。前四行回答不同来源怎样产生候选监督信号，后两行回答样本标注怎样共同使用、检查和修正。来源与后续处理承担不同职责，不是六种互斥算法，也不要求每个方案全部采用。}
  \label{tab:knowledge-acquisition-chapter-map}
\end{table*}

\begin{table*}[htbp]
  \centering
  \small
  \setlength{\tabcolsep}{3pt}
  \begin{tabularx}{\linewidth}{@{}>{\RaggedRight\setlength{\parindent}{0pt}\arraybackslash}p{15mm}>{\RaggedRight\setlength{\parindent}{0pt}\arraybackslash}p{37mm}>{\RaggedRight\setlength{\parindent}{0pt}\arraybackslash}p{28mm}>{\RaggedRight\setlength{\parindent}{0pt}\arraybackslash}p{43mm}>{\RaggedRight\setlength{\parindent}{0pt}\arraybackslash}X@{}}
    \toprule
    \textbf{来源} & \textbf{准备与投入} & \textbf{典型标注结果} & \textbf{限制与易错处} & \textbf{核验入口} \\
    \midrule
    人工 & 规范、界面与培训；逐项判断和复核 & 类别、区域、解释或弃权 & 受能力与可见证据限制；歧义、疲劳和偏好可致错 & 独立专家、预先确认的参考答案与留出抽查 \\
    规则 & 名录、编号映射与适用条件；开发和版本维护 & 类别、关联与约束 & 编号误配、未收录对象和名录过期 & 错配反例、名录版本核对与覆盖抽查 \\
    模型 & 模型与预测请求的组织；调用、筛选和更新 & 类别、区域、关系或生成材料 & 依赖能力与分布；高置信错误和来源偏差可被复制 & 独立标注、可验证规则与任务留出评价样本集 \\
    样本本身 & 观测与构造依据；样本处理、构造变换和训练 & 恢复目标、视图关系与分组 & 依赖保留的结构；错误配对、改变物种含义的变换，或依赖背景完成训练 & 检查构造、移除相应目标作对照、独立下游评价 \\
    \bottomrule
  \end{tabularx}
  \caption{在同一花卉任务下选择监督来源。各行并列列出准备与逐项投入、典型标注结果、容易出错的条件和核验办法；依据这些条件选择来源，而不能仅凭标注结果形式或自动化程度判断。弃权是人工来源的动作，不是额外的类别答案。}
  \label{tab:knowledge-acquisition-source-conditions}
\end{table*}

\subsection{多来源结果的对齐与融合}
\label{sec:knowledge-acquisition-alignment-fusion}

单独比较来源以后，可以把互补的办法放进同一个稀有花卉方案。植物学家先为少量代表照片作答并写出规范，规则按采集名录取得更多候选类别，视觉模型为尚无名录答案的照片提出预测，同一对象的多角度照片则构造关系目标。它们分别补充已有缺口，因而可以同时采用，而不是必须挑一个来源处理全部问题。

共同使用以前，先确认不同结果回答的是同一个问题。人工对照片甲中的左侧花朵作答，模型对整张照片作答，两份结果的对象不同；人工用“鸢尾”指一个具体物种，规则用它指一组近似类别，两份结果的含义也不同。对象对齐（object alignment）把结果连到同一对象，语义对齐（semantic alignment）确认类别名称、任务和适用范围采用同一含义。两种对齐完成之后，才能把不同来源的“鸢尾”作为可比较的票。

多份可比较判断仍可能冲突，因此需要融合（fusion），即按照明确办法把多来源信息共同用于形成结果。一个最小办法是对同一对象的类别票按多数决定类别；更复杂的办法还会估计来源在不同条件下的可靠性。多数票也有前提：如果三条规则都从同一份错误名录复制答案，它们的三票只是同一错误的重复，不能当作三份独立支持。恢复目标和关系目标则没有与物种类别相同的回答空间，可以分别进入训练，无须融合成一个物种名称。

已有模型还可以作为整体共同使用，例如让几个模型各自预测后再组合输出，或者由一个模型向另一个模型提供训练目标。这与综合同一照片的标签票是不同层面的操作。第~\ref{chap:knowledge-fusion}章会分别处理多来源标注与已有模型的联合使用；在当前获取方案中，需要先说明交付的是标注内容还是模型，以及后续怎样使用它。

\subsection{获取、检查与再获取}
\label{sec:knowledge-acquisition-feedback}

把这些操作组织起来，一次性方案可以从试点进入批量获取，再按需要对齐融合、抽查修正并验收交付。如果抽查发现稀有物种在遮挡下经常被误判，持续方案还要把这个结果送回下一轮，改选样本、追加人工判断或调整来源与预算。每轮维护日志应说明改变了什么，哪些旧结果的条件仍成立、可以保留，哪些对象必须重新获取，以及何时停止。只是再次调用同一个模型生成同类答案，没有补入独立标注依据时，增加轮次并不能作为可用监督信号增加的证据。

例如，试点发现大多数错误来自看不见花瓣的照片，下一轮可以只对这些对象补拍局部，再交给专家判断；若错误来自旧名录中的名称映射，则应更新映射并重查受影响的记录。前一种处理增加可见证据，后一种处理改变答案解释依据。检查结果之所以形成反馈（feedback），是因为它决定了下一轮改什么，而不仅是把同样的获取步骤再执行一次。

图~\ref{fig:knowledge-acquisition-workflow}把样本标注的去向与方案的调整分开画出。四类来源都能送入候选集合，实线监督信号路径随后分叉：多份可比较结果需要对齐融合，单一来源或独立训练目标可以直接检查。虚线反馈则从质量判断返回来源选择或获取机制，表示检查结果改变下一轮怎样做。沿这条反馈看，发现遮挡下的错误可能只要求重做特定对象，而不是把全部步骤无条件重跑。

\begin{figure*}[htbp]
  \centering
  \begin{tikzpicture}[
  x=1mm,y=1mm,
  every node/.style={font=\figurefont\footnotesize,text=BookInk},
  stage/.style={rectangle,draw=BookRule,line width=0.8pt,fill=BookPaper,text width=24mm,minimum height=13mm,align=center,inner sep=2mm},
  source/.style={rectangle,rounded corners=1.2mm,draw=FigureInput,line width=0.8pt,fill=FigureInput!8,text width=14mm,minimum height=10mm,align=center,inner sep=1.5mm},
  fusion/.style={stage,draw=FigureModel,fill=FigureModel!8,text width=20mm},
  check/.style={stage,draw=BookAmber,fill=BookAmber!8,text width=21mm},
  result/.style={stage,draw=FigureOutput,fill=FigureOutput!6,text width=15mm},
  flow/.style={knowledge arrow,draw=BookTeal,line width=0.9pt,rounded corners=1.2mm},
  control/.style={knowledge arrow,draw=BookMuted,line width=0.8pt,rounded corners=1.2mm},
  bus/.style={draw=BookMuted,line width=0.8pt},
  knowledge bus/.style={draw=BookTeal,line width=0.9pt},
  feedback bus/.style={draw=BookAmber,dashed,line width=0.8pt,rounded corners=1.2mm},
  adjust/.style={knowledge arrow,draw=BookAmber,dashed,line width=0.8pt,rounded corners=1.2mm},
  edge label/.style={font=\figurefont\scriptsize,inner sep=1pt,text=BookMuted}
]
  \node[stage] (need) at (15,12) {需求与资源\\范围、效果、预算};
  \node[stage] (select) at (47,12) {来源选择\\对象、机制、投入};

  \node[source] (human) at (83,36) {人工};
  \node[source] (rule) at (83,20) {规则};
  \node[source] (model) at (83,4) {模型};
  \node[source] (data) at (83,-12) {样本构造};
  \node[font=\figurefont\scriptsize\bfseries,text=BookTeal] at (83,46) {执行选定的获取机制};

  \node[stage,draw=FigureModel,fill=FigureModel!8] (candidate) at (126,12) {候选监督信号\\含来源与状态};
  \node[fusion,text width=24mm] (align) at (83,-46) {对象与语义对齐\\需要时再融合};
  \node[check] (quality) at (126,-46) {质量判断\\覆盖与可靠性\\下游效用};
  \node[result] (usable) at (159,-46) {可用结果\\或停止};

  \draw[control] (need.east) -- (select.west);
  \draw[control] (select.east) -- (67,12);
  \draw[bus] (67,-12) -- (67,36);
  \draw[control] (67,36) -- (human.west);
  \draw[control] (67,20) -- (rule.west);
  \draw[control] (67,4) -- (model.west);
  \draw[control] (67,-12) -- (data.west);

  \draw[knowledge bus] (100,-12) -- (100,36);
  \draw[flow] (human.east) -- (100,36);
  \draw[flow] (rule.east) -- (100,20);
  \draw[flow] (model.east) -- (100,4);
  \draw[flow] (data.east) -- (100,-12);
  \draw[flow] (100,12) -- (candidate.west);

  \draw[flow] ([xshift=-8mm]candidate.south) -- (118,-25) -| (align.north);
  \node[edge label,anchor=south] at (99,-24) {多份可比结果};
  \draw[flow] (align.east) -- (quality.west);
  \draw[flow] (candidate.south) -- (quality.north);
  \node[edge label,anchor=west,align=left] at (129,-16) {单一来源\\或独立目标};
  \draw[flow] (quality.east) -- (usable.west);

  \coordinate (feedback split) at (100,-64);
  \draw[feedback bus] (quality.south) |- (feedback split);
  \draw[adjust] (feedback split) -- (100,-72) -| (select.south);
  \node[edge label,text=BookAmber,anchor=south] at (70,-71) {调整来源选择};
  \draw[adjust] (feedback split) -- (64,-64) -- (64,-23) -- (67,-23) -- (67,-12);
  \node[edge label,text=BookAmber,anchor=south] at (82,-63) {调整获取机制};

  \draw[flow] (4,-82) -- (16,-82);
  \node[anchor=west,font=\figurefont\scriptsize] at (18,-82) {监督信号流动};
  \draw[adjust] (57,-82) -- (69,-82);
  \node[anchor=west,font=\figurefont\scriptsize] at (71,-82) {质量反馈};
\end{tikzpicture}

  \caption{获取、融合与质量改进的一种组织方式。实线监督信号路径由四类来源汇入候选集合，按是否需要综合判断选择融合或直接检查；虚线表示依据覆盖、可靠性与效用调整来源选择和获取机制。需求与来源选择之间的连接组织执行，不传递新的标签或关系。图不要求每次启用全部来源或全部环节。}
  \label{fig:knowledge-acquisition-workflow}
\end{figure*}

比较两个完整方案时，固定的不只是花卉分类这个任务名称，还包括目标地区与时间范围、评价样本、训练设置和成本核算周期。覆盖、可靠性、下游指标与分项投入需要同时报告，准备和持续成本还要按来源列明。如果根据同一评价样本集反复选样、调整规则或改写模型提示，这组评价样本就已经参与方案选择，不能再当作独立的最终效果证据，否则评价可能偏向被选中的方案。跨来源预算协调见第~\ref{chap:knowledge-fusion}章，人工内部的选样、选人与重复标注分配见第~\ref{chap:knowledge-acquisition-human}章；日志、存储与调度的系统实现不在本部分展开。

\section{知识获取的研究历史}
\label{sec:knowledge-acquisition-history}

前节的方案既可以逐例取得答案，也可以把专家判断写成规则反复执行。历史中同样能看到这两种安排。Rosenblatt在1958年的感知机研究中讨论用带答案样本调整连接；感知机根据输入特征的加权结果给出类别，样本答案用于判断是否需要调整权重。1984年出版的MYCIN项目回顾则说明怎样把领域专家的判断依据编码为条件规则，并检查、维护规则库。前者让答案随样本进入学习，后者把判断依据写成可重复执行的规则。两者为理解获取方式提供早期例证，并不是监督学习或专家系统的唯一起源\scite{rosenblatt1958,buchanan1984rulebased}

逐例提供答案需要持续投入人工，规则库需要维护覆盖范围与例外。后续研究从这些不同限制出发：有限答案应当用在哪些对象上，多份不一致的答案怎样使用，没有任务标签时还能从哪里取得学习目标，以及模型生成内容后怎样检查和保留它们。本节按这些问题组织代表工作，说明每项工作改变的机制。自动化让人参与的环节发生变化，但界定任务、处理难例和核验结果仍可能需要人工，这些路线也继续并存。

\subsection{有限标签的选择与利用}
\label{sec:knowledge-acquisition-history-label-budget}

当未标注照片很多、预算只允许追加少量人工答案时，我们需要选择询问对象。设现有已标注照片恰好都是“紫色鸢尾”和“红色玫瑰”，两个候选预测办法可能都答对：一个主要依赖颜色，另一个主要依赖花瓣形态。它们都符合现有答案，却可能在一张粉色鸢尾照片上给出不同预测。询问这种照片，有机会知道哪种判断办法需要修正；继续询问两者都答对的近重复照片，未必能提供同样的区分信息。这是一个教学情形，用来说明有限标签为什么可能留下多个不同候选。

Seung、Opper与Sompolinsky在1992年的委员会查询（query by committee）中，利用这种预测分歧选择对象。委员会是一组当前允许考虑的候选假设，每个假设是一种从输入得到类别的预测办法；让它们对未标注样本分别作答，再优先询问分歧大的对象。取得人工答案以后，就能据此更新候选及其判断。查询价值取决于允许的假设范围、委员会能否反映当前不确定性，以及人能否给出有依据的答案。若分歧只是因为照片严重模糊、任何人都无法辨认，花费同样预算未必取得有效区分\scite{seung1992querycommittee}

另一条路线利用已有少量标签和大量未标注样本，减少新答案必须逐个来自人工的程度。Blum与Mitchell在1998年的协同训练（co-training）中，让两个分类器相互提供候选标签。它们观察同一对象的两种特征视图（feature view），即两组不同的可用特征。原工作中的网页例子有助于理解：网页正文是一组描述，指向该网页的链接文字是另一组描述，两者都可能提供网页类别的线索。

先用少量已标注网页分别训练两个分类器，再让它们给未标注网页预测。一个分类器选出的高置信候选标签，可以成为另一个分类器后续训练的依据；两者使用的特征不同，因而有机会互相补入信息。原工作的理论设定要求每个视图都足以预测类别，并要求给定类别后两视图条件独立：在已知类别的条件下，再知道正文视图，不会改变链接视图的条件分布，反之亦然。仅仅把特征分成两组并不能满足这些条件。两视图共享偏差时，可能反复把同一种错误传给对方。委员会查询主要改变向人询问的对象，协同训练改变标签怎样在模型之间扩展；扩展出来的仍是候选答案\scite{blum1998cotraining}

选样还有一个起步问题：标签极少时，新训练的任务模型尚未形成可靠判断，用它的预测分数来选择待标注样本也可能不可靠。这称为冷启动（cold start）困难，即缺少足够初始依据，后续选择机制难以开始有效工作。一个思路是借助已在其他材料上训练的大语言模型，也就是能处理和生成文本的已有大规模模型。2026年发表于TACL的ActiveLLM在文本少样本场景中利用大语言模型选择待人工标注的实例。它提供询问优先级，最终答案仍需人工取得；效果证据限于论文的模型、样本集与选择协议，不能直接推广到所有任务\scite{bayer2026activellm}

\subsection{多来源答案的质量推断}
\label{sec:knowledge-acquisition-history-source-quality}

一张照片交给三个人，两人判断玫瑰、一人判断鸢尾，多数票会选择玫瑰，但票数本身没有说明三人的专业能力和照片难度。某人平时可靠，也可能在遮挡照片上出错；另一人的少数意见则可能来自更充分的专业知识。Whitehill等人在2009年的模型中，联合推断图像标签、标注者能力和样本难度：用同一个概率模型描述不同能力的人在不同难度样本上怎样产生答案，再依据多次作答估计这些未知量。这样，来源与样本条件进入了综合判断，票的作用不必始终相同。推断仍依赖模型设定和作答信息，不能仅凭某人总与多数一致就证明他正确\scite{whitehill2009crowdlabels}

专家还可以把经验写成程序，让它反复提交候选票。Ratner等人在2017年发布的Snorkel用标注函数（labeling function）实现这一安排：程序检查指定条件，满足时返回候选类别，没有类别判断时返回不投票编码。例如，前节按采集编号查名录的规则就可以写成这样的函数。多条函数可能覆盖不同照片，也可能在同一照片上冲突，因而还需要估计哪些函数较可靠、哪些票共享依据。

Snorkel使用生成式来源模型（generative source model）描述这些票怎样产生。这里建模的对象是来源输出：假定存在尚未确定的真实类别，再用来源准确性和相关结构解释为什么会看到这一组类别票。观察多条函数在许多对象上的同意与分歧，帮助估计来源性质，随后形成用于训练的综合标签。它减少了逐例人工作答，却把一部分工作转为编写函数、规定来源模型和检查其条件。若多条函数复制同一错误，或者观测不足以区分来源质量与真实类别，综合结果仍可能不可靠。未触发与主动弃权都可能被程序编码为“不投票”，原始触发结果与来源动作仍需按前节约定分别保存\scite{ratner2017}

\subsection{从原内容构造学习目标}
\label{sec:knowledge-acquisition-history-internal-target}

如果观测很多，却没有目标任务的类别答案，仍可以利用原内容构造学习要求。前节以图像遮挡说明了这一点，文本也有相应机制。设原句是“花瓣呈紫色”，我们遮住“紫色”，让模型根据其余内容预测原来被遮住的单位；在这个教学例子中，原句自身就提供了一个可检查的答案，无需另外询问句子的分类标签。文本先切分成词元（token），即模型读取和预测的基本单位，实际词元可以是词，也可以是词的一部分。

BERT采用的遮蔽语言模型（masked language model）目标，要求根据上下文预测选中位置的原词元。它通过大量这样的预测学习文本表示。先用一般学习目标训练模型，作为后续具体任务的起点，称为\term{预训练}（pre-training）。取得的直接依据是原词元恢复目标，训练的结果是模型及其表示；原文本并没有因此自动得到情感类别或其他下游标签。BERT于2018年发布预印本、2019年发表于NAACL，在这里代表内容恢复路线，并不是自监督学习的唯一起点\scite{devlin2019}

从一般目标学得的表示到具体任务输出，还缺少一层联系。一个模型能根据上下文恢复词元，不等于它已经知道当前情感任务有哪些类别；同样，图像模型能恢复花瓣，也不等于它已经会输出本地区的物种名称。任务适配（task adaptation）就是把已有模型接到具体任务的输入与输出要求上，例如利用已有表示，再用物种标签训练分类输出。若继续调整已有模型的参数，这种后续训练称为\term{微调}（fine-tuning）。预训练提供起点，适配确定怎样用于当前任务，任务评价再检查结果。

Bommasani等人在2021年的\term{基础模型}（foundation model）报告中，把在大规模且来源多样的训练样本上训练、随后可适配多种下游任务的模型作为共同研究对象。一个已训练模型可以服务多种任务，准备投入因而有机会被复用；报告也强调下游系统会继承上游模型的缺陷。回到花卉例子，构造恢复目标、学得适合识别的表示、取得正确物种答案是三种不同结果，前一种成立不能代替后两种的检查\scite{bommasani2021foundation}

\subsection{模型生成材料与递归使用}
\label{sec:knowledge-acquisition-history-generated-material}

模型来源还可以从“给已有样本作答”扩展到“连问题和材料一起生成”。Self-Instruct让语言模型生成指令、输入材料和示范答案，再过滤无效或过于相似的项目，用于后续微调。指令（instruction）说明要执行什么任务，输入材料说明这一次对什么执行，示范答案则给出训练要模仿的输出。例如，“提取句中的花卉名称”是任务指令，“花坛里种着鸢尾”是这次输入，“鸢尾”是示范答案。三者共同组成一项训练实例，任何一项不匹配都可能产生错误监督。这个例子只说明组成方式，不代表论文使用的具体实例。

相比预测一张已有照片的物种，这种安排扩大了可以构造的训练材料，也扩大了需要检查的对象：指令是否明确、输入是否包含足够依据、答案是否符合指令和材料，以及新任务是否只是在重复少数模式。过滤能排除指定的不合格项，却不能保证留下的答案事实正确或任务分布完整。Self-Instruct提供的是生成与筛选训练实例的办法，生成数量不能单独作为可靠知识增加的证据\scite{wang2023selfinstruct}

生成候选以后，还可以按明确标准筛选。DeepSeek-R1的2025年原始报告（arXiv v1）组合了模型生成、拒绝采样、可验证任务的规则奖励与一般场景的奖励模型。这里的拒绝采样（rejection sampling）指报告所用的候选筛选环节：先生成，再拒绝不合格项、保留合格项。奖励（reward）则为训练结果提供评分反馈；规则奖励按可执行标准评分，奖励模型（reward model）通过另一个模型计算分数。以可核对最终答案的数学题为例，检查规则能判断最终答案是否符合标准，却不自动证明每一步解释都正确。一般任务中，模型评分也可能继承评分模型的偏差。获取、筛选和反馈可以相连，但每种检查都要说明它能覆盖哪些属性；本章不展开强化学习的训练算法\scite{deepseek2025r1}

模型提供标签时，还会遇到来源能力与学习者能力不一致的问题。弱到强泛化（weak-to-strong generalization）研究把弱模型生成的标签用于训练更强模型，考察学习者能否取得超过标签来源的任务结果。这里的强弱指同一任务设置下相对的模型能力。强学习者已有的训练基础可能比弱监督者更丰富，弱标签只是后续指导它的一部分信息；它并不是从零开始、只能照抄这批标签。不过，弱标签中的系统错误仍可能被学习者继承。该工作于2023年公开预印本、2024年发表于ICML，在若干任务上观察到学习者超过弱监督者，但仍未恢复强模型的全部能力。结果支持特定设置下弱监督可能有价值，不能作为弱来源普遍足以监督更强系统的保证\scite{burns2023weaktostrong}

使用一次模型生成的训练实例，与反复依赖前代模型生成材料，是不同的训练安排。合成训练样本（synthetic training sample）指由模型等生成机制构造、用于训练的实例。递归使用时，第1代模型产生第2代的训练材料，第2代再产生第3代的材料，后续观测越来越依赖前代模型保留了什么。若生成器很少产生稀有鸢尾，下一代就可能更少见到这类样本；每轮总条数不变，也可能逐渐只剩常见玫瑰。这是用于说明覆盖退化的花卉情形，不是论文中的花卉实验。

Shumailov等人在递归合成训练的理论与实验设置中研究\term{模型坍塌}（model collapse）：原分布中少见的情形逐代丢失，后代模型越来越难保留这些尾部内容。结论依赖具体的递归训练机制，不能推出任意一次使用合成样本都会退化。它要求我们在方案中记录来源谱系（source lineage），即样本由哪个模型、依据哪一代材料生成，以及独立于这条生成链的原始观测是否仍被保留。只统计每代条数，无法检查稀有内容是否已经丢失\scite{shumailov2024collapse}

递归训练还要区分替换与累计。替换会丢弃旧训练集、只保留新合成样本；累计则保留最初从真实对象取得的样本，并继续加入历代合成样本。以教学计数为例，最初有1000张真实照片，每代生成200张，第5轮累计方案保留$1000+5\times200=2000$张，替换方案只保留最近200张。累计至少没有按这项规则主动丢弃最初的照片，同时增加了存储与重训投入；这个条数计算没有证明识别能力更好。

Gerstgrasser等比较这两种安排，其线性回归分析和有限实验表明，累计可以避免所考察替换机制中的误差持续累积，但结论依赖抽样、噪声与模型条件。筛选办法、旧样本权重或目标分布改变以后，还须重新检查稀有内容的覆盖。保留真实样本是方案的一个机制条件，不能被简化为“只要累计就能保证任意递归训练有效”\scite{ka-gerstgrasser2024-accumulation}

这些问题可以再放回年份顺序中对照。表~\ref{tab:knowledge-acquisition-history-works}按所引版本的年份排列代表工作，用相同栏目比较研究问题、获取机制、新增依据与验证边界。它还列出后续主讲章节，便于从当前问题继续阅读。各行的任务和证据不同，不能合并性能数字作统一排名；发表较晚也不说明在当前花卉任务上更适合。

\begin{table*}[p]
  \centering
  \footnotesize
  \setlength{\tabcolsep}{2.5pt}
  \begin{tabularx}{\linewidth}{@{}>{\RaggedRight\setlength{\parindent}{0pt}\arraybackslash}p{28mm}>{\RaggedRight\setlength{\parindent}{0pt}\arraybackslash}p{29mm}>{\RaggedRight\setlength{\parindent}{0pt}\arraybackslash}p{34mm}>{\RaggedRight\setlength{\parindent}{0pt}\arraybackslash}p{34mm}>{\RaggedRight\setlength{\parindent}{0pt}\arraybackslash}X@{}}
    \toprule
    \textbf{代表工作} & \textbf{研究问题} & \textbf{来源或获取机制} & \textbf{新增标签、关系或训练目标} & \textbf{主要假设、验证边界与后续入口} \\
    \midrule
    Rosenblatt感知机（1958） & 怎样依据带答案样本调整模型参数 & 人工答案随样本进入模型参数更新 & 把观测与类别答案组成可学习样本 & 只代表早期监督学习实例；模型准备见第~\ref{chap:knowledge-acquisition-model-based}章 \\
    MYCIN回顾（1984） & 怎样获取并维护专家判断依据 & 访谈专家并把专家判断编码为规则库 & 可重复执行、检查和修订的规则说明 & 规则覆盖与例外依赖领域维护；见第~\ref{chap:rule-based-knowledge-acquisition}章 \\
    委员会查询（1992） & 哪些未标注样本值得询问 & 候选假设的预测分歧指导人工查询 & 把有限人工答案集中到当前模型难以确定的样本 & 依赖假设空间与委员会；人工选择见第~\ref{chap:knowledge-acquisition-human}章 \\
    协同训练（1998） & 怎样利用少量标签和大量未标注样本 & 两种特征视图的分类器相互提供候选标签 & 扩展可训练的样本标签 & 依赖视图充分性与给定类别后的条件独立性；样本关系见第~\ref{chap:data-derived-knowledge-acquisition}章 \\
    Whitehill等（2009） & 怎样融合能力不同的标注者 & 联合建模标签、标注者能力和样本难度 & 估计综合标签及来源可靠性 & 观测能否区分真值与来源质量、模型设定是否合适，均需检验；人员与融合见第~\ref{chap:knowledge-acquisition-human}章和第~\ref{chap:knowledge-fusion}章 \\
    Snorkel（2017） & 怎样反复执行并组合专家启发式 & 返回类别票或不投票编码的标注函数与来源模型 & 批量候选标签及来源准确性、依赖估计 & 规则覆盖与相关结构可能设错；规则与融合见第~\ref{chap:rule-based-knowledge-acquisition}章和第~\ref{chap:knowledge-fusion}章 \\
    BERT（2018/2019） & 怎样让无下游标签文本参与预训练 & 遮蔽观测并预测原词元 & 内容恢复目标和预训练表示 & 内部目标不等于任务答案；样本构造见第~\ref{chap:data-derived-knowledge-acquisition}章 \\
    基础模型报告（2021） & 怎样理解大规模预训练与任务适配 & 在多种来源的训练样本上训练后适配多种任务 & 对可复用表示与共同问题的研究框架 & 下游会继承上游缺陷；模型来源见第~\ref{chap:knowledge-acquisition-model-based}章 \\
    Self-Instruct（2023） & 怎样扩大指令训练材料 & 模型生成指令、预测请求和示范答案并过滤 & 新的任务描述与训练实例 & 过滤不保证事实正确或分布完整；模型来源见第~\ref{chap:knowledge-acquisition-model-based}章 \\
    弱到强泛化（2023预印本／2024发表） & 弱来源能否支持更强学习者 & 弱模型标签监督强模型 & 特定实验中的弱监督训练信号 & 未恢复全部强模型能力；质量边界见第~\ref{chap:knowledge-quality-assessment-correction}章 \\
    递归合成训练样本研究（2024） & 模型生成的训练样本反复回流会发生什么 & 前代模型生成后代训练样本 & 对覆盖退化条件的证据 & 结论依赖递归合成样本的训练机制；质量与修正见第~\ref{chap:knowledge-quality-assessment-correction}章 \\
    DeepSeek-R1原始报告（2025） & 怎样组合生成样本与筛选反馈 & 模型生成、拒绝采样、规则及奖励模型 & 推理与非推理训练材料 & 验证器和生成模型都有覆盖边界；见第~\ref{chap:knowledge-acquisition-model-based}章和第~\ref{chap:knowledge-quality-assessment-correction}章 \\
    ActiveLLM（2026） & 少样本冷启动时怎样选择待标注实例 & 大语言模型根据任务描述选择样本 & 面向人工查询的样本优先级 & 依赖模型和选择协议；人工选样见第~\ref{chap:knowledge-acquisition-human}章 \\
    \bottomrule
  \end{tabularx}
  \caption{按获取问题横向比较代表研究。各行分别说明机制、新增标签、关系或训练目标及其使用边界，并给出后续入口，不比较跨任务性能。BERT和弱到强泛化同时列预印本与正式发表年份，其余采用所引版本年份；代表工作不等于相应方向的唯一起点。}
  \label{tab:knowledge-acquisition-history-works}
\end{table*}

横向比较以后，图~\ref{fig:knowledge-acquisition-history}再显示这些问题在哪些年份同时发展。四条泳道按选择人工判断、组织多来源与规则、利用样本内在结构、模型生成与核验来归位，同一工作也可能服务于不止一条路线。横向位置保留年份先后，但间距不按时间长度缩放；泳道与连接说明本章怎样组织问题，不证明直接继承或因果关系。人工标注和专家规则放在持续背景中，不为这些实践强行指定一个统一起点。

\begin{figure*}[!t]
  \centering
  \begin{tikzpicture}[
  x=1mm,y=1mm,
  every node/.style={font=\figurefont\scriptsize,text=BookInk},
  lane label/.style={anchor=west,font=\figurefont\bfseries\scriptsize,text width=20mm,align=left},
  timeline/.style={draw=BookRule,line width=0.8pt},
  event/.style={rectangle,rounded corners=1mm,line width=0.8pt,minimum height=9mm,align=center,inner sep=1.2mm},
  selection event/.style={event,draw=FigureInput,fill=FigureInput!8,text width=20mm},
  source event/.style={event,draw=FigureModel,fill=FigureModel!8,text width=20mm},
  data event/.style={event,draw=BookTeal,fill=BookTeal!7,text width=20mm},
  generation event/.style={event,draw=BookAmber,fill=BookAmber!8,text width=20mm},
  stem/.style={draw=BookMuted,line width=0.6pt},
  time arrow/.style={knowledge arrow,draw=BookMuted,line width=0.8pt}
]
  \fill[BookPaper] (0,-82) rectangle (22,43);
  \draw[BookRule,line width=0.6pt] (22,-82) -- (22,43);

  \node[lane label,text=FigureInput] at (1,28) {选择人工判断};
  \node[lane label,text=FigureModel] at (1,9) {组织多来源与规则};
  \node[lane label,text=BookTeal] at (1,-10) {利用样本内在结构};
  \node[lane label,text=BookAmber] at (1,-31) {模型生成与核验};

  \node[anchor=south west,font=\figurefont\scriptsize\bfseries,text=BookTeal] at (24,46) {持续背景：人工逐例标注、专家规则、任务设计与核验};
  \draw[time arrow] (24,42) -- (166,42) node[above left,font=\figurefont\scriptsize] {年份：位置保持先后，间距不按比例};
  \draw[timeline] (24,28) -- (166,28);
  \draw[timeline] (24,9) -- (166,9);
  \draw[timeline] (24,-10) -- (166,-10);
  \draw[timeline] (24,-31) -- (166,-31);

  \node[selection event] (qbc) at (32,28) {1992\\委员会查询};
  \node[selection event,text width=17mm] (activellm) at (159,28) {2026\\ActiveLLM};

  \node[source event,text width=22mm] (whitehill) at (64,9) {2009\\Whitehill等};
  \node[source event] (snorkel) at (88,9) {2017\\Snorkel};

  \node[data event] (cotraining) at (49,-10) {1998\\协同训练};
  \node[data event,text width=22mm] (bert) at (95,-10) {2018/2019\\BERT};
  \node[data event,text width=22mm] (foundation) at (120,-10) {2021\\基础模型报告};

  \foreach \yearx/\eventyear in {131/2023,142/2024,153/2025}{
    \fill[BookAmber] (\yearx,-31) circle (1.2pt);
    \node[anchor=south,font=\figurefont\scriptsize] at (\yearx,-29) {\eventyear};
  }
  \node[generation event] (self instruct) at (102,-43) {2023\\Self-Instruct};
  \node[generation event,text width=28mm] (weak strong) at (102,-66) {2023预印本／2024发表\\弱到强泛化};
  \node[generation event,text width=22mm] (collapse) at (151,-49) {2024\\递归生成退化};
  \node[generation event,text width=22mm] (deepseek) at (151,-75) {2025\\DeepSeek-R1};
  \draw[stem] (131,-31) -- (131,-43) -- (self instruct.east);
  \draw[stem] (131,-43) -- (131,-66) -- (weak strong.east);
  \draw[stem] (142,-31) -- (142,-39) -| (collapse.north);
  \draw[stem] (153,-31) -- (166,-31) |- (deepseek.east);

\end{tikzpicture}

  \caption{获取问题随时间扩展而并存。四条泳道按研究问题归位，横向位置保留年份先后，间距不按比例；下方引线把生成相关工作接到年份点，不表示方法继承。图不表达因果或后来路线取代早期路线，“2018/2019”分别指BERT预印本与正式发表。}
  \label{fig:knowledge-acquisition-history}
\end{figure*}

\FloatBarrier
\section{本章小结}
\label{sec:knowledge-acquisition-summary}

回到开篇的花卉模型，照片提供观测，物种标签提供检查类别预测的答案，损失函数再把预测与答案的差异变成训练依据。没有物种名称时，对象关系仍能要求表示接近，类别约束能排除不合法判断，恢复目标也能从观测中构造学习要求。它们都可以指导训练，但给出的信息不同，恢复得好或满足一个约束，都不能直接代替正确的物种答案。

获取这些内容以后，还要保存对象、来源、内容、适用范围和状态。执行成功只说明操作完成，内容已给出只说明取得判断，按指定依据核验后才有相应的质量结论。未知、弃权和不适用之所以分开，是因为它们会带来补证据、恢复执行、转交或排除等不同处理。这样保存，后续训练就能分辨一条负标签与一条尚无答案的记录。

选择获取方案时，需要从实际缺口出发，再决定来源和处理方式。人工、规则、模型与样本本身提供的依据和易错条件不同，可以互补；可比较答案要先对齐对象与含义，不同训练目标则可以分别使用。记录能否读取、参考证据是否支持、用于训练是否帮助目标任务，应分别检查。在预算、可靠性和覆盖满足要求的方案中，再比较任务效用；候选生成数、置信度或训练损失都不能单独替代这些判断。

代表性研究表明，获取已经从逐例答案和专家规则扩展到查询选择、来源建模、内部目标以及模型生成与筛选，但这些路线并未相互取代。批量处理带来新的获取机会，也需要应对来源共享错误、候选回流和覆盖退化；生成或筛选能力扩大，仍不保证新增样本标注可靠、有用。后续各章据此展开四类来源的具体机制，再处理融合与质量改进。本章确定了需求、样本标注和评价口径，具体算法怎样实现、在什么设置下值得采用，还须沿这些入口继续检验。
